\documentclass{article}
\usepackage[UTF8]{ctex}
\usepackage{amsmath,amssymb,amsfonts}
\usepackage{enumitem}
\usepackage{geometry}
\geometry{a4paper, margin=1in}
% === 错题本排版优化（不参与题目标识） ===
\setlength{\parindent}{0pt}
\setlength{\parskip}{0.35em plus 0.10em minus 0.08em}
\setlength{\emergencystretch}{3em}
\setlist{itemsep=0.25em, topsep=0.40em, parsep=0pt}
\allowdisplaybreaks[2]
% === 错题本排版优化结束 ===
\begin{document}
\section*{安徽大学2026年数学分析真题}
\section*{112．安徽大学2026年数学分析试卷}

\subsection{题目1}
1．设 \(x_{1}=\frac{1}{5}, x_{n}>0, n=1,2, \cdots\) ，且成立等式 \(x_{n+1}^{2}-5 x_{n}-6=0\) ，证明：数列 \(\left\{x_{n}\right\}\) 收敛，并求其极限．

\subsection{题目2}
2．计算下列极限：

\begin{itemize}
  \item[（1）] \(\lim _{x \rightarrow 0} \frac{\sqrt{1+4 x}-2 \ln (x+1)-1}{e^{x}-\sin 3 x+2 x-1}\) ．
  \item[（2）] \(\lim _{n \rightarrow \infty} \sum_{k=1}^{n} \frac{\sin \frac{\pi}{n}}{2+\cos \frac{k \pi}{n}}\) ．
\end{itemize}

\subsection{题目3}
3．设广义积分 \(\int_{a}^{+\infty} f(x) \mathrm{d} x\) 收敛，且导函数 \(f^{\prime}(x)\) 在 \([a,+\infty)\) 上有界，证明： \(\lim _{x \rightarrow+\infty} f(x)=0\) ．

\subsection{题目4}
4．设 \(f(x)\) 在 \([a, b]\) 上具有连续导数，且 \(f(a)=0\) ，证明：

\begin{itemize}
  \item[（1）] \(\int_{a}^{b} f^{2}(x) \mathrm{d} x \leq \frac{(b-a)^{2}}{2} \int_{a}^{b}\left[f^{\prime}(x)\right]^{2} \mathrm{~d} x\) ．
  \item[（2）] \(\int_{0}^{2 \pi} f(a \cos x+b \sin x) \mathrm{d} x=\int_{0}^{2 \pi} f\left(\sqrt{a^{2}+b^{2}} \sin x\right) \mathrm{d} x\) ．
\end{itemize}

\subsection{题目5}
5．设函数 \(f(x)\) 在区间 \((-1,1)\) 上有二阶连续导数，且 \(\lim _{x \rightarrow 0} \frac{f(x)}{x}=0, f^{\prime \prime}(x) \neq 0\) ，证明：级数\(\sum_{n=2}^{\infty} f(\sqrt{\sqrt[n]{2}-1}\) 发散。

\subsection{题目6}
6．设 \(u_{n}(x)(n=1,2,3, \cdots)\) 在 \([a, b]\) 上单调，且函数项级数 \(\sum_{n=1}^{\infty} u_{n}(x)\) 在 \(x=a\) 和 \(x=b\) 处收敛。

\begin{itemize}
  \item[（1）] 证明：\(\sum_{n=1}^{\infty} u_{n}(x)\) 在 \([a, b]\) 上一致收敛。
  \item[（2）] 证明：\(\sum_{n=2}^{\infty} \ln \left(1+\frac{x}{n \ln ^{3} n}\right)\) 在 \([-a, a]\) 上一致收敛，其中 \(0<a<2 \ln ^{3} 2\) ．
\end{itemize}

\subsection{题目7}
7．设一元函数 \(y=\varphi(x)\) 在 \(x=0\) 附近有定义，\(\varphi(0)=0\) ，且在 \(x=0\) 附近满足 \(|\varphi(x)| \leq x^{2}\) ，定义二元函数

\[
f(x, y)= \begin{cases}\varphi\left(\frac{x y}{\sqrt{x^{2}+y^{2}}} \ln \left(1+\sqrt{x^{2}+y^{2}}\right)\right), & (x, y) \neq(0,0) \\ 0, & (x, y)=(0,0)\end{cases}
\]

判断 \(f(x, y)\) 在点 \((0,0)\) 处是否可微，并证明．

\subsection{题目8}
8．计算曲线积分

\[
\oint_{L}\left(\frac{-y}{x^{2}+y^{2}}+2 y\right) \mathrm{d} x+\left(\frac{x}{x^{2}+y^{2}}+3 y\right) \mathrm{d} y .
\]

其中曲线 \(L\) 的方程为 \(\frac{x^{2}}{9}+\frac{y^{2}}{4}=1\) ，方向为逆时针方向。

\subsection{题目9}
9．第一卦限中，在曲面 \(\frac{x^{2}}{a^{2}}+\frac{y^{2}}{b^{2}}+\frac{z^{2}}{c^{2}}=1\) 上求一点使得该点切平面与三个坐标平面所围成的立体体积最小。

\subsection{题目10}
10．设区域 \(\Omega\) 由分片光滑的封闭曲面 \(\sum\) 围成，\(u(x, y, z)\) 在 \(\bar{\Omega}\) 上具有二阶连续偏导数，且在 \(\bar{\Omega}\) 上满足

\[
\frac{\partial^{2} u}{\partial x^{2}}+\frac{\partial^{2} u}{\partial y^{2}}+\frac{\partial^{2} u}{\partial z^{2}}=0
\]

（1）证明： \(\iint_{\sum} \frac{\partial u}{\partial n} \mathrm{~d} S=0\) ，其中 \(\mathbf{n}\) 是曲面 \(\sum\) 的单位外法向量。\\
（2）设 \(\left(x_{0}, y_{0}, z_{0}\right) \in \Omega\) 为一定点，证明：

\[
u\left(x_{0}, y_{0}, z_{0}\right)=\frac{1}{4 \pi} \iint_{\sum}\left(u \frac{\cos (\widehat{r, n})}{r^{2}}+\frac{1}{r} \frac{\partial u}{\partial n}\right) \mathrm{d} S
\]

其中 \(r=\left(x-x_{0}, y-y_{0}, z-z_{0}\right), r=\sqrt{\left(x-x_{0}\right)^{2}+\left(y-y_{0}\right)^{2}+\left(z-z_{0}\right)^{2}}\) ．

\subsection{题目11}
11．计算积分 \(\int_{0}^{\pi} \ln \left(1+\frac{1}{2} \cos x\right) \mathrm{d} x\) ．
\end{document}
